problem:  
Let \( (M^n,g) \) be a compact, simply connected Riemannian manifold with positive sectional curvature \(K>0\).  
At each point \(p\in M\), denote the curvature operator  
\[
R_p:\Lambda^2T_pM\to\Lambda^2T_pM
\]
with eigenvalues \(0<\lambda_1(p)\le\cdots\le\lambda_N(p)\), \(N=\tfrac{n(n-1)}{2}\).  
Define, for \(a\in(0,1)\), the *traceless‑spectral cone*
\[
\mathcal{C}_a=\{R>0: \tfrac{\lambda_1(p)}{\operatorname{tr}(R_p)}\ge a\}.
\]
Note that the normalization by the trace makes \(\mathcal{C}_a\) an \(\mathrm{O}(n)\)-invariant, scale‑free subset of curvature operators, representing pinching toward constant curvature along the line spanned by \(\mathrm{Id}_{\Lambda^2}\).

**Problem:**  
For each \(n\ge4\), determine whether there exists \(a_n>0\) such that the cone \(\mathcal{C}_{a_n}\) is preserved under the Hamilton ODE  
\[
\partial_tR = R^2+R^\#,
\]
and furthermore, that if a compact \(M^n\) has curvature operators in \(\mathcal{C}_{a_n}\) at time \(0\), then the normalized Ricci flow exists for all time and converges to a constant‑curvature metric.  
Equivalently, does a uniform lower bound on the minimum eigenvalue relative to the trace guarantee convergence of the curvature operator to the constant‑curvature ray under the flow?  
If not, can one characterize the maximal open convex subset of positive curvature operators that both (i) contain the constant‑curvature ray and (ii) are invariant under the ODE?

Why is it a "good" problem:  
This formulation turns the spectral pinching idea into an **invariance classification** problem for the Hamilton ODE, aligning directly with the modern program pioneered by Hamilton, Böhm–Wilking, and Brendle–Schoen. It asks for a quantitative, *spectrally defined* invariant cone bridging curvature algebra and flow dynamics: a natural intermediary between known cones (positive, 2‑positive, PIC) and full constant curvature.  

The problem probes how far the Ricci‑flow invariance of curvature positivity can be relaxed toward *near‑isotropy*, defined not by specific quadratic inequalities but by spectral proximity to the identity. Discovering or ruling out an invariant threshold \(a_n\) would clarify whether spectral pinching alone—without explicit algebraic curvature identities—suffices to control Ricci‑flow evolution.  

It is expressible in elementary terms yet conceptually deep, asking for the *border between algebraic positivity and dynamical stability* of curvature under flow. Progress would enhance our understanding of invariant curvature cones and possibly reveal new geometric evolution principles. Hence, it stands as a meaningful and forward‑looking research problem at the geometry and analysis interface.